% =============================================================================
%  题名页、摘要、译者说明、目录
% =============================================================================
\title{\bfseries 局部同调镜像对称\\[6pt]
       \large Local Homological Mirror Symmetry}
\author{Si-Yang Liu\thanks{伦敦帝国理工学院数学系，伦敦 SW7 2AZ。
        电子邮箱：\texttt{[邮箱隐去]}}
        \and Kyungmin Rho\thanks{波恩马克斯·普朗克数学研究所。
        电子邮箱：\texttt{[邮箱隐去]}}}
\date{}

\maketitle

\begin{center}
\small
中译注本\quad 底本：arXiv:2611.07202v1 [math.SG]（2026 年 11 月 4 日，共 56 页）\\[2pt]
译者：数译（DSH 数学文献翻译工作组）\\
橙色「译注」框、紫色「概念回顾」框、蓝色「本节重点」框均为译者所加，原文没有。
\end{center}

\vspace{0.8em}

\section*{摘要}
\addcontentsline{toc}{section}{摘要}
局部镜像对称源于 Chiang--Klemm--Yau--Zaslow 的工作，它提示了向量丛的镜像与其零截面的镜像之间可能存在联系。本文为这类局部簇构造了同调意义下的镜像，并证明了相应的同调镜像对称陈述。作为应用，我们为 Hirzebruch 面构造了无穷多个异域 Weinstein 扇形，作为它们的镜像。

\vspace{0.6em}
\noindent\textbf{关键词：}同调镜像对称；局部镜像对称；包裹 Fukaya 范畴；等变层；Künneth 公式

\clearpage

\section*{译者说明}
\addcontentsline{toc}{section}{译者说明}

\begin{enumerate}[label=（\arabic*）,leftmargin=2.4em,itemsep=0.35em]
  \item \textbf{底本。}译文依据 arXiv:2611.07202v1 [math.SG]（2026 年 11 月 4 日，56 页）的 LaTeX 源文件与排版 PDF 逐段对照译出。原文的节、小节、定理、公式、图表编号一律保留原编号；译文正文中每页起始处以小字注明原文页码，以便与原文逐页核对。
  \item \textbf{定理与公式编号。}原文所有定理类环境（定理、定义、命题、推论、问题、例、引理、注、约定、练习）共用一个计数器、按节编号，因此「定义 3.5」之后可能是「注 3.6」。中译本保持这一编号规则。原文带编号的公式只有几处，其余公式均不编号，译文亦同。
  \item \textbf{术语。}全书术语以项目冻结术语表为准，附于附录 B；术语首次出现时以「中文（English）」形式给出英文，其后只用中文。Fukaya、Liouville、Weinstein、Künneth、Floer、Čech、Orlov、Hirzebruch、Calabi--Yau 等专名与已通用的英文术语按惯例保留原文。人名音译，首次出现附原文拼写。
  \item \textbf{三类方框。}译文含三种原文没有的方框，均可由导言区开关一键关闭：
    \begin{itemize}[leftmargin=1.2em,itemsep=0.1em]
      \item 橙色\textbf{译注}框（全书连续编号）：易错预警、含糊点拨、背景直觉；
      \item 紫色\textbf{概念回顾}框：对非专家不显然的概念作 3--6 句回顾，附 nLab 或原始文献出处；
      \item 蓝色\textbf{本节重点}框：每节末尾给出该节在全文证明链条中的位置。
    \end{itemize}
    三者都只是「译者所加」的阅读辅助，不改变正文内容；正文与方框请分开阅读。
  \item \textbf{原文笔误。}原文的明显笔误（拼写、重复词、个别上下标）在正文中照原样排印，并在译注中说明；凡属译者判断而非原文明确写出的内容，一律放入译注或概念框，不混入正文。
  \item \textbf{图表。}原文 5 张图与全部交换图（tikz-cd）按原样重排，图题译为中文；图中英文标签在保持数学含义的前提下译出。
  \item \textbf{编译。}源码用 XeLaTeX 连续编译两遍；参考文献沿用原文的 \texttt{alphaurl} 样式，条目保留英文，引用标签与原文一致。
\end{enumerate}

\vspace{0.5em}
\noindent{\small 译者按：原文第 1 节末尾含作者关于人工智能使用情况的说明，译文照译，见第 1 节末。}

\clearpage
\tableofcontents
\clearpage
